\chapter{为什么超对称性可能重要}

\lectureoverview{真空涨落将本章的各个主题联系起来。边界条件产生卡西米尔力，圈效应威胁希格斯质量，新的虚粒子改变跑动的规范耦合，而真空能支配宇宙膨胀。这些问题促使人们考虑低能超对称性，以此稳定希格斯能标、提供暗物质候选者并改善规范耦合统一。随后我们使用格拉斯曼坐标构造超场，并利用别列津积分提取超对称作用量。}

\section{反射板之间的真空涨落}

我们回到超对称性，首先解决一个关于卡西米尔效应的问题。将两块平行的反射板放置在相距\(d\)处。即使不存在任何实粒子，量子场也具有真空涨落。我们可以把它们想象为虚过程，或者等价地，想象为场的振子模式的零点运动。板施加了边界条件，因此它们之间允许的涨落不同于无限制空间中的涨落。

重要的量不是能量的绝对零点，而是间距变化时所产生的变化。将单位面积的真空能记为\(\mathcal E(d)=E(d)/A\)。一旦\(\mathcal E\)依赖于\(d\)，它就是势能，而单位面积上的力为
\begin{equation}
\frac{F}{A}=-\frac{d\mathcal E}{dd}.
\end{equation}
因此，虚过程的重排变成可观测的力。在熟悉的理想导体板电磁学例子中，能量和压强为
\begin{align}
\mathcal E(d)&=-\frac{\pi^2}{720\,d^3},\\
\frac{F}{A}&=-\frac{\pi^2}{240\,d^4},
\end{align}
单位为\(\hbar=c=1\)。\editorialnote{这些系数以及\(d^{-3}\)能量每面积与\(d^{-4}\)压强之间的区别，解决了黑板上提出的量纲问题。它们是标准的理想板结果，而不是在此进行的推导。} 负号表示吸引。

\begin{classroomqa}
\audiencequestion{四次方描述的是力还是能量？ \lecturetimestamp{00:04:39}}
\lecturerresponse{对于四维时空中的平行板，单位面积的能量按\(d^{-3}\)下降，而对\(d\)求导使单位面积的力按\(d^{-4}\)下降。一旦恢复\(\hbar\)和\(c\)，量纲分析就已经确定了这些幂次。}
\end{classroomqa}

人们有时把该计算总结为板中断了虚光子圈。只要不把它理解得太字面化，这一图像是有用的。精确的计算是允许的场模式谱的变化，或者等价地，是量子有效作用量的变化。无论用哪种语言，都会得到相同的依赖于间距的能量。

\begin{classroomqa}
\audiencequestion{卡西米尔力可能是排斥的而不是吸引的吗？ \lecturetimestamp{00:07:48}}
\lecturerresponse{可以。符号并非普适的；它取决于场、几何、材料和边界条件。费米子圈通常带来与玻色子圈相反的圈符号，这促使了黑板上提出的建议，但费米子统计本身并不是一条普遍定理，断言每个卡西米尔力都必须为排斥。\editorialnote{这一限定取代了课堂上故意试探性的回答，代之以避免错误普适规则所需的一般陈述。}}
\end{classroomqa}

\subsection{什么可以作为真空能的证据？}

卡西米尔力已被测量，原子可观测量的确也对真空圈敏感。然而，这两个事实本身都不能证明绝对真空能会产生引力。卡西米尔实验测量的是由边界引起的能量差；同一个力也可以在不赋予无界绝对零点能求和以操作意义的情况下被描述。

决定性的检验来自引力。能量是引力荷。如果空空间携带均匀的能量密度，它就必须进入爱因斯坦方程并影响宇宙膨胀。因此，宇宙的加速膨胀就是真空样成分的直接引力证据。

\begin{classroomqa}
\audiencequestion{为什么通俗介绍把真空能称为神秘的，仿佛物理学家不知道它是什么？ \lecturetimestamp{00:08:19}}
\lecturerresponse{量子场论赋予每个场一个零点贡献，因此真空能的可能性并不神秘。神秘之处是定量上的：为什么其引力值比基于已知量子场和高能截断的估计小得如此非凡？令人惊讶的事实是它几乎不存在，而不是它仅仅存在。}
\end{classroomqa}

多年来，观测给出的上限如此之小，以至于在心理上很容易让人把宇宙学常数精确设为零，并去寻找某种强制其相消的原理。改进后的宇宙学观测却反而发现了一个小的非零值。这一发现并没有使真空能在概念上变得陌生；它使那个无法解释的压低程度变得不可忽视。相关的比较取决于所假定的紫外标度，但每一个自然的高能估计都要求巨大的相消。

\section{一个超荷对多粒子态的作用}

在列举相信超对称的理由之前，我们应当先消除一个可能的悖论。玻色子可以任意大量地占据一个模式，而一个费米子模式的占据数为零或一。如果一种对称性交换玻色子和费米子，它会不会把一个模式中的一千个玻色子变成一千个被禁戒的费米子？

不会。一个费米型生成元一次只改变一个量子。在振子模型中，设\(a\)湮灭一个玻色子，\(c^\dagger\)产生其费米子伙伴。一个超荷可以表示为
\begin{equation}
Q=c^\dagger a.
\end{equation}
对于一个有\(n\)个玻色子且没有费米子的态，
\begin{equation}
Q\ket{n,0}
=\sqrt n\,\ket{n-1,1}.
\end{equation}
一个玻色子被移除，一个费米子被插入。再次用同一个荷作用则给出零：
\begin{equation}
Q^2=(c^\dagger)^2a^2=0.
\end{equation}
这就是从费米型产生算符继承来的幂零性。

\begin{classroomqa}
\audiencequestion{如果许多玻色子占据同一个态，为什么一次超对称变换不会在那个态中产生许多费米子并违反泡利不相容原理？\lecturetimestamp{00:12:05}}
\lecturerresponse{生成元并不会替换整个占据数。它使态改变一个单位的费米子数。同一个费米型荷不能像普通的转动生成元那样被复合，因为它的平方为零。一个互补的荷可以撤销这次交换，而混合反对易子则闭合到时间平移上。}
\end{classroomqa}

最后一句至关重要。简写\(Q^2=0\)适用于单个手征分量，或适用于刚刚写出的那个简单荷；它并不意味着超荷的每一个乘积都为零。在四维超对称中，
\begin{equation}
\{Q_\alpha,\bar Q_{\dot\beta}\}
=2\sigma^\mu_{\alpha\dot\beta}P_\mu,
\end{equation}
而
\begin{equation}
\{Q_\alpha,Q_\beta\}=0,
\qquad
\{\bar Q_{\dot\alpha},\bar Q_{\dot\beta}\}=0
\end{equation}
对于没有中心荷的最小代数成立。\editorialnote{这种分量形式精确地表达了课堂上的简写：同一个\(Q\)的重复作用为零，而一个超荷与其共轭则生成平移。}

\section{寻找低能超对称的三个理由}

\begin{classroomqa}
\audiencequestion{除了相消真空能之外，超对称还有别的动机吗？\lecturetimestamp{00:15:06}}
\lecturerresponse{有三个密切相关的动机：控制希格斯自能、可能存在稳定的暗物质粒子，以及规范耦合在高能下更精确地相交。它们都不能证明超对称在自然界中得以实现，但合在一起使超对称成为一个值得检验的严肃理论。}
\end{classroomqa}

\emph{低能超对称}这个说法是相对的。它并不意味着超伴子在日常尺度上很轻；它意味着它们的质量并不远高于电弱标度和加速器标度。如果把超伴子推到任意高的能标，这三个论证都会减弱。

\subsection{保护希格斯标度}

标量质量对高能物理异常敏感。示意性地，一个玻色子圈图产生如下形式的修正
\begin{equation}
\delta m_H^2\big|_{\mathrm B}
\sim +\frac{\lambda}{16\pi^2}\Lambda^2,
\end{equation}
其中\(\Lambda\)表示有效理论被使用到的标度。一个费米子圈图带有相反的符号，
\begin{equation}
\delta m_H^2\big|_{\mathrm F}
\sim -\frac{|y|^2}{8\pi^2}\Lambda^2.
\end{equation}
系数取决于相互作用和场的多重数，但超对称将它们联系起来，使得主导的紫外敏感性在伙伴之间相消。\editorialnote{这些示意性的圈图公式揭示了定性自能论证背后的数学。它们并不打算作为完整的标准模型希格斯质量计算。}

这一点很重要，因为希格斯期望值设定了电弱标度。规范玻色子和费米子的质量服从诸如
\begin{equation}
m_W=\frac{gv}{2},
\qquad
m_f=\frac{y_fv}{\sqrt2}.
\end{equation}
这样的关系。如果希格斯部门被自然地推到巨大的标度，那么除非裸参数以极高的精度针对辐射修正进行调整，否则熟悉的低能粒子谱就会丧失。

精确的超对称会使相消精确成立，但精确的超对称已经被观测到的玻色子和费米子质量不相等所排除。在软破缺超对称下，二次敏感性被替换为对超伴子质量劈裂的依赖，示意性地
\begin{equation}
\delta m_H^2
\sim
\frac{g^2}{16\pi^2}
\left(m_{\mathrm{soft}}^2-m_{\mathrm{ordinary}}^2\right)
\ln\frac{\Lambda^2}{m_{\mathrm{soft}}^2}.
\end{equation}
只有当劈裂不太大时，这种相消才仍然有用。

\subsection{暗物质级联的一个稳定终点}

将谱分为普通粒子及其超对称伙伴。费米子有玻色型伙伴；玻色子有费米型伙伴。如果相互作用保持一种区分这两类的宇称，那么一个超对称伙伴的衰变必须留下另一个超对称伙伴：
\begin{equation}
\widetilde X
\longrightarrow X+\widetilde Y
\longrightarrow X+Y+\widetilde Z
\longrightarrow\cdots.
\end{equation}
级联过程不能丢弃最后一个超对称伙伴。因此它终止于最轻的超对称粒子，即LSP。

常用的对称性是\(R\)宇称，
\begin{equation}
R_p=(-1)^{3(B-L)+2s}.
\end{equation}
标准模型粒子具有\(R_p=+1\)，而它们的超对称伙伴具有\(R_p=-1\)。如果\(R_p\)是精确的，那么每个相互作用都包含偶数个超对称伙伴，且LSP是稳定的。\editorialnote{仅靠超对称性并不能保证LSP稳定。稳定性要求守恒的\(R\)宇称或具有相同结果的其他对称性；没有它，LSP可能衰变。}

暗物质必须至少存活宇宙学时间尺度。它还必须足够弱地相互作用，以便反复穿过可见星系盘而不损失大量能量。带电荷或强相互作用的粒子群会散射、耗散能量，并倾向于像普通物质一样坍缩到同一薄层结构中。相反，观测到的延展晕指向一种中性、弱相互作用的成分。

这并不能唯一地确定LSP。取决于模型和对称性破缺模式，候选者可能是电弱规范微子和希格斯微子的中性混合态，也可能是引力微子。带电荷或带色荷的候选者受到强烈限制。超对称性所提供的并不是一个保证的答案，而是一种产生中性、长寿命终点的自然机制。

\section{跑动耦合与统一的迹象}

第三个动机从电荷屏蔽出发。测得的耦合依赖于能标，因为虚粒子会修饰相互作用。对于电磁学，插入光子传播子中的电子-正电子圈在长距离上屏蔽电荷。在更短距离，或更高动量转移下，这种极化云位于探针之间的部分更少，因此有效电荷更大。

写出
\begin{equation}
\alpha_i=\frac{g_i^2}{4\pi}.
\end{equation}
在一圈阶，逆耦合随能量对数线性跑动：
\begin{equation}
\frac{d\alpha_i^{-1}}{d\ln\mu}
=-\frac{b_i}{2\pi}.
\end{equation}
因此
\begin{equation}
\alpha_i^{-1}(\mu)
=\alpha_i^{-1}(\mu_0)
-\frac{b_i}{2\pi}\ln\frac{\mu}{\mu_0}.
\end{equation}
\editorialnote{这个单圈方程是作为对数能量函数绘制的近乎直线图像的精确版本。阈值和更高圈修正会使这些线略微弯曲。}

在电弱能标以上，三个基本规范耦合分别是
\begin{equation}
U(1)_Y,\qquad SU(2)_L,\qquad SU(3)_c.
\end{equation}
的耦合。低能下用于直觉的电磁耦合是电弱对称性破缺后前两者的组合。对于QCD，非阿贝尔规范场产生反屏蔽：强耦合在短距离下减小，即渐近自由现象。因此，\(\alpha_3^{-1}\)随\(\ln\mu\)上升，而其他逆耦合具有不同的斜率。

标准模型的曲线暗示性地彼此指向，但不在单一点相交。超对称伙伴修改虚圈，从而改变系数\(b_i\)。在传统的大统一归一化下，
\begin{align}
(b_1,b_2,b_3)_{\mathrm{SM}}
&=\left(\frac{41}{10},-\frac{19}{6},-7\right),\\
(b_1,b_2,b_3)_{\mathrm{MSSM}}
&=\left(\frac{33}{5},1,-3\right).
\end{align}
当超对称伙伴阈值接近TeV能标时，外推的超对称曲线在
\begin{equation}
M_{\mathrm{GUT}}\sim 10^{16}\ \mathrm{GeV}.
\end{equation}
附近以高得多的精度相交。
没有任何实验能直接达到该能标。低能耦合被测量，场内容决定跑动，曲线被外推许多个数量级。两条线总会在某处相交；三条独立曲线在一点附近相交才是非平凡的观察。

\begin{classroomqa}
\audiencequestion{改进的外推是否假设了精确超对称性和相等的伙伴质量？\lecturetimestamp{00:38:09}}
\lecturerresponse{不。超对称性是破缺的，因此质量是分裂的。重要的是超对称伙伴圈何时进入跑动。如果伙伴不是太重，斜率会足够早地改变，使三个耦合收敛。如果它们极其重，超对称阈值到来得太晚，改进的相交就会丧失。}
\end{classroomqa}

阈值修正和未知的高能谱阻止了三条线以数学精确性相交的字面主张。尽管如此，结果仍然引人注目：同一个有助于希格斯并提供了热暗物质候选者的近似超对称伙伴能标，也改善了统一。

\subsection{引力与共同的高能标度}

引力之所以表现不同，是因为它的荷就是能量本身。在质心能量\(E\)处，一个有用的无量纲量度是
\begin{equation}
\alpha_G(E)\sim GE^2
=\left(\frac{E}{M_{\mathrm{Pl}}}\right)^2.
\end{equation}
在普通能量下，这个量小得不可思议，但它随\(E\)二次增长。因此它的倒数下降得比按对数跑动的规范耦合倒数快得多，并在普朗克标度附近变为1阶。

这一比较暗示，所有相互作用可能由相邻的高标度支配，但这不如规范统一那样干净。规范耦合是可重整场论中的无量纲参数；牛顿常数是有量纲的，因此要把引力放到同一张图上，需要约定一种约定。

\begin{classroomqa}
\audiencequestion{超对称伙伴是否参与同样的三种规范相互作用，还是它们定义了新的力？\lecturetimestamp{00:42:35}}
\lecturerresponse{它们携带与其伙伴相同的规范量子数。一个超电子具有电子的电荷，一个超夸克具有对应夸克的色荷和电弱荷，尽管它们的自旋和质量不同。因此它们的圈图修正的是同样的规范相互作用，而不是引入三种新的相互作用。}
\end{classroomqa}

\begin{classroomqa}
\audiencequestion{在LHC发现一个超对称伙伴，是否会成为超对称的决定性证据？\lecturetimestamp{00:43:27}}
\lecturerresponse{仅凭一个陌生粒子的名称本身并不足够，但一个遵循自旋、电荷、质量和耦合之间超对称关系的超对称伙伴谱将极具特征性。超对称施加许多相互关联的限制，因此详细的相互作用模式将提供决定性的检验。}
\end{classroomqa}

\subsection{为什么标度不能无代价地被移开}

如果超对称伙伴只稍微超出实验可及范围，那么这三种动机仍然都能起作用。如果它们的质量被提高几个数量级，每个论据都会变差：
\begin{enumerate}
\item 希格斯质量抵消会留下对分裂的很大残余敏感度。
\item 常规弱标度热遗迹不再在预期标度上出现。
\item 规范跑动阈值上移，三个耦合相交的精度变差。
\end{enumerate}
超对称仍可能存在于高得多的标度上，但它将不再以同样经济的方式解决这些特定的低能问题。一个想法寿命长并不是自然界必须采用它的证据；而是表明该提议足够自洽，值得进行决定性检验。

\section{宇宙学能告诉我们什么、不能告诉我们什么关于暗物质}

\begin{classroomqa}
\audiencequestion{标准宇宙学能否像核合成约束氢和氦那样，确定暗物质的总量以及每个暗物质粒子的质量？\lecturetimestamp{00:46:32}}
\lecturerresponse{宇宙学能很好地测量暗物质的总质量密度，但仅凭密度无法把质量与数目分开。同样的密度可以由少数重粒子携带，也可以由数量极多的极轻粒子携带。要识别微观粒子，还需要额外的理论或非引力观测。}
\end{classroomqa}

两种宽泛的可能性说明了这种歧义。一个接近电弱标度或TeV标度的弱相互作用粒子可以从热浴中冻结出来。在另一个极端，一个极轻的玻色子，例如轴子，可以占据一个相干场位形或凝聚体。中间可能性并不会仅仅因为两个发展最充分的机制位于相反两端附近就被禁止。

晕结构本身提供了一个重要约束。暗物质穿过普通物质时几乎没有摩擦，因此它保持空间延展，而不会坍缩进发光盘。旋转曲线和其他引力现象告诉我们质量存在，但它们本身并不清点粒子数目。

\begin{classroomqa}
\audiencequestion{早期宇宙最初是否可能只包含超粒子，后来它们衰变成我们看到的粒子？\lecturetimestamp{00:49:11}}
\lecturerresponse{在远高于所有相关质量的温度下，快速反应驱动各物种趋向热平衡，而不是任意的一边倒布居。普通碰撞在超对称伙伴不足时产生它们，而超对称伙伴碰撞和衰变则补充普通粒子。精确丰度取决于相互作用和自由度，但热等离子体不断趋向一种平衡。}
\end{classroomqa}

随着宇宙冷却，有质量的粒子种类受到玻尔兹曼抑制，并发生湮灭或衰变。稳定的LSP在其相互作用速率变得慢于宇宙膨胀时，就不再跟踪平衡态：
\begin{equation}
\Gamma_{\mathrm{ann}}
\sim n\langle\sigma v\rangle
\lesssim H.
\end{equation}
剩余的粒子形成遗迹丰度。\editorialnote{这一冻结判据是热平衡论证的标准定量完成。其他LSP候选者，尤其是耦合极弱的引力微子，可以具有非热历史。}

\begin{classroomqa}
\audiencequestion{LSP是哪种普通粒子的伙伴，它必须不相互作用吗？\lecturetimestamp{00:51:05}}
\lecturerresponse{不存在与模型无关的答案。一个可行的暗物质LSP预期是电中性且无色的，但它不必完全不相互作用。它可以是与光子、弱规范玻色子和希格斯场相关的混合态，也可以是引力微子。它的确切身份由超对称破缺谱决定。}
\end{classroomqa}

\section{真空图、自能图与膨胀}

\begin{classroomqa}
\audiencequestion{希格斯发散与真空能发散之间有什么区别？\lecturetimestamp{00:52:09}}
\lecturerresponse{真空图没有外粒子线；它改变空空间的有效作用量，并对真空能作出贡献。希格斯自能图有两条外希格斯腿；它改变希格斯传播子，从而改变其质量。两者都涉及虚涨落，但它们回答的是不同的问题。}
\end{classroomqa}

人们可以把一个实粒子想象成插入真空中的杂质。附着在这个杂质上的真空涨落会改变它的传播和质量。卡西米尔板也改变真空涨落，但通过宏观边界条件。共同的要素是量子涨落对所插入之物的响应；在每种情形中被修正的可观测量不同。

\vspace{-2pt}
\subsection{物质、辐射与真空的稀释方式不同}

设\(a(t)\)为宇宙标度因子，
\begin{equation}
H=\frac{\dot a}{a}
\end{equation}
为哈勃参数。均匀宇宙中的局部应力-能量守恒给出
\begin{equation}
\dot\rho+3H(\rho+p)=0.
\end{equation}
对于状态方程\(p=w\rho\)，积分得到
\begin{equation}
\rho(a)\propto a^{-3(1+w)}.
\end{equation}
\editorialnote{连续性方程提供了粒子计数解释和波长拉伸解释背后的紧凑推导。}

对于冷的非相对论性物质，\(p\approx0\)，因此
\begin{equation}
\rho_m\propto a^{-3}.
\end{equation}
固定的共动体积按\(a^3\)增长，而其粒子所携带的静止能量近似保持不变。

对于辐射，\(p=\rho/3\)，并且
\begin{equation}
\rho_r\propto a^{-4}.
\end{equation}
三个幂次来自数密度的稀释。第四个幂次来自红移：
\begin{equation}
\lambda\propto a,
\qquad
E_\gamma=\frac{2\pi\hbar c}{\lambda}\propto a^{-1}.
\end{equation}

\begin{classroomqa}
\audiencequestion{如果光子有波函数，那么在推导稀释定律时，为什么我们可以把它说成粒子而不是波？\lecturetimestamp{00:58:21}}
\lecturerresponse{这一区分不影响标度关系。在粒子描述中，共动体积内的数目固定，每个量子按\(a^{-1}\)红移。在经典波描述中，波长拉伸，场能量下降相同的量。我们可以使用任何一种使膨胀最容易看清的语言。}
\end{classroomqa}

真空能具有\(p_\Lambda=-\rho_\Lambda\)，因此
\begin{equation}
\rho_\Lambda=\text{constant}.
\end{equation}
它不随空间膨胀而稀释。确实，一个共动区域包含
\begin{equation}
E_\Lambda=\rho_\Lambda a^3V_0,
\end{equation}
它随其物理体积增长。这一事实产生了我们稍后将讨论的能量守恒问题。

不同的幂次组织了宇宙历史。辐射下降最快，物质下降较慢，而真空能不下降。因此，一个微小的真空项在早期宇宙中可以忽略不计，而在晚期占主导。使用这里讨论的粗略数值，当前能量预算约为\(70\%\)真空能和\(30\%\)物质，辐射则小得多。

\section{几何即能量}

爱因斯坦方程的时间-时间分量把前面的标度律转化为动力学。板书从一个口号和一个几何项开始。

\begin{figure}[H]
\centering
\includegraphics[width=0.68\textwidth]{lecture_06_figure_03.png}
\lectureframe{01:06:37.088}
\caption{组织性口号\emph{几何即能量}，其下写有膨胀项\(H^2=(\dot a/a)^2\)。}
\end{figure}

对于均匀且各向同性的宇宙，第一弗里德曼方程为
\begin{equation}
H^2+\frac{k}{a^2}
=\frac{8\pi G}{3}\rho,
\end{equation}
在单位\(c=1\)下，其中\(k=+1,0,-1\)表示闭合、平坦或开放的空间几何，并且
\begin{equation}
\rho
=\rho_{\mathrm{ordinary}}
+\rho_{\mathrm{dark}}
+\rho_{\mathrm{radiation}}
+\rho_{\mathrm{vacuum}}.
\end{equation}
在\(1/a^2\)旁看到的符号取决于曲率项是保留在左边还是移到右边，也取决于\(k\)的值。

\begin{figure}[H]
\centering
\includegraphics[width=0.96\textwidth]{lecture_06_figure_04.png}
\lectureframe{01:08:46.794}
\caption{完整的平衡关系：几何一侧是膨胀与空间曲率，源一侧是物质、辐射与真空能。物质又进一步分为普通物质与暗物质两部分。}
\end{figure}

\subsection{宇宙学哈密顿约束}

\begin{classroomqa}
\audiencequestion{如果共动区域中的真空能随宇宙膨胀而增长，那么是否存在一个恒定的净能量密度，或者一个守恒的总能量？\lecturetimestamp{01:01:51}}
\lecturerresponse{广义相对论并不为任意膨胀时空提供一个普适的全局能量。然而，在均匀宇宙学约化中，时间-时间爱因斯坦方程是一个哈密顿约束。当把几何与物质放在同一侧时，受约束的哈密顿量为零。变化的物质项由引力自由度来平衡，而不是由盒中一个普通的固定能量来平衡。}
\end{classroomqa}

用符号表示，这一正则表述为
\begin{equation}
\mathcal H_{\mathrm{gravity}}
+\mathcal H_{\mathrm{matter}}
=0.
\end{equation}
这就是“总能量为零”这一说法的受控含义。它是由时间重参数化不变性所强加的运动方程，而不是断言每个局域能量密度都为零。在共动体积中，普通物质的能量大致固定，辐射能量减少，真空能增长；与此同时，引力贡献与膨胀率作出调整，使爱因斯坦约束仍然得到满足。

局域守恒仍然严格成立：
\begin{equation}
\nabla_\mu T^{\mu\nu}=0.
\end{equation}
对于均匀流体，这恰好就是连续性方程
\(\dot\rho+3H(\rho+p)=0\)。一般而言失效的是平直时空的直觉，即认为在膨胀宇宙中还必须存在一个与坐标无关的、守恒的全局物质能量总和。\editorialnote{这一区分使“能量始终为零”这一课堂说法变得精确：局域协变守恒与宇宙学哈密顿约束是严格成立的，而一个唯一的全局物质加引力能量一般并没有定义。}

\begin{classroomqa}
\audiencequestion{零约束是决定宇宙是开放、闭合还是平坦的条件吗？\lecturetimestamp{01:04:42}}
\lecturerresponse{不是。哈密顿约束在每一种情形下都成立。开放、平坦和闭合宇宙通过曲率标签\(k\)来区分，因而也通过\(k/a^2\)项来区分。曲率改变的是平衡关系中的一个贡献；它并不会开启或关闭该约束。}
\end{classroomqa}

\section{格拉斯曼坐标与超场}

\begin{classroomqa}
\audiencequestion{格拉斯曼数真的出现在理论的拉格朗日量中吗？\lecturetimestamp{01:09:37}}
\lecturerresponse{它们首先出现，是因为费米子场是格拉斯曼奇的。它们第二次出现，是作为超空间的坐标。超场依赖于这些奇坐标，而超对称作用量则是通过对它们求导和积分来构造的。}
\end{classroomqa}

普通场是时空的函数，
\begin{equation}
\phi=\phi(x^\mu).
\end{equation}
四维最小超对称把\(x^\mu\)扩充为一个二分量奇旋量\(\theta^\alpha\)及其共轭\(\bar\theta^{\dot\alpha}\)：
\begin{equation}
z^M=(x^\mu,\theta^\alpha,\bar\theta^{\dot\alpha}).
\end{equation}
这里有两个复奇分量，或者说四个实奇方向。超场是这个超空间上的函数：
\begin{equation}
\Phi=\Phi(x,\theta,\bar\theta).
\end{equation}

这些并不是度规意义上的微小空间维度。格拉斯曼坐标没有正的数值大小。称它小，只是对幂零性的一种助记：
\begin{equation}
(\theta^\alpha)^2=0,
\qquad
\theta^\alpha\theta^\beta=-\theta^\beta\theta^\alpha.
\end{equation}
\editorialnote{幂零性，而非数值上的小，才是任何在有限多个格拉斯曼坐标中的展开都会终止的原因。}

有四个独立的奇坐标时，一个完全一般的超场至多有\(2^4=16\)个格拉斯曼单项式。用符号表示，
\begin{equation}
\begin{aligned}
\Phi(x,\Theta)={}&\phi(x)+\Theta_A\psi_A(x)\\
&+\frac{1}{2}\Theta_A\Theta_B Z_{AB}(x)+\cdots\\
&+\Theta_1\Theta_2\Theta_3\Theta_4D(x),
\end{aligned}
\end{equation}
其中\(\Theta_A\)统称四个实奇坐标。由于整个超场可以选为偶的，偶次幂的系数是玻色性的，奇次幂的系数是费米性的。因此，一个超场就打包了一个由普通分量场组成的有限多重态。

这种打包不只是记号。超对称在几何上表现为奇方向上的平移与时空平移的混合。按照超空间规则由超场构造出的拉格朗日量，会以与由张量构造出的标量拉格朗日量继承旋转或洛伦兹不变性相同的方式继承该对称性。

作用量现在包含格拉斯曼测度。一个典型的全超空间项具有形式
\begin{equation}
S_D=\int d^4x\,d^2\theta\,d^2\bar\theta\,
\mathcal K(\Phi,\Phi^\dagger).
\end{equation}
手征项则使用
\begin{equation}
S_F=\int d^4x\,d^2\theta\,W(\Phi)+\mathrm{h.c.}
\end{equation}
\editorialnote{全超空间测度与手征测度之间的区分，是对示意性\(d^4x\,d^4\theta\)作用量的标准细化。具体例子将决定哪种测度是合适的。}

因此，我们遇到了一个普通微积分无法回答的问题：对一个反交换坐标的积分意味着什么？

\section{Berezin积分}

\subsection{普通模型：全导数}

从一个普通函数\(F(x)\)出发，它在\(x\)的正负大值处都趋于零。微积分基本定理给出
\begin{equation}
\int_{-\infty}^{\infty}dx\,\frac{dF}{dx}
=F(\infty)-F(-\infty)
=0.
\end{equation}

\begin{figure}[H]
\centering
\includegraphics[width=0.88\textwidth]{lecture_06_figure_05.png}
\lectureframe{01:26:31.594}
\caption{Berezin积分的普通原型：全导数的积分归结为端点值，当函数衰减时为零。黑板下沿部分被遮挡，因此正文中给出了清晰的方程。}
\end{figure}

我们希望Grassmann积分是线性的，并且具有全导数积分为零这一有用性质。对于一个奇坐标，要求
\begin{align}
\int d\theta\,[af(\theta)+bg(\theta)]
&=
a\int d\theta\,f(\theta)
+b\int d\theta\,g(\theta),\\
\int d\theta\,\partial_\theta f(\theta)&=0.
\end{align}
第二个条件是测度平移不变性的代数类比。奇数的数轴没有端点，也没有对微小区间的求和。

\begin{classroomqa}
\audiencequestion{如果没有不定Grassmann积分，我们是否要放弃积分是微分之逆这一想法？\lecturetimestamp{01:21:19}}
\lecturerresponse{是的。Berezin积分是一个确定的线性泛函，而不是求原函数的运算。对于一个Grassmann变量，它实际上由与微分相同的代数运算表示。}
\end{classroomqa}

\begin{classroomqa}
\audiencequestion{对全时空的\(d^4x\)积分也是一个定积分吗？\lecturetimestamp{01:21:34}}
\lecturerresponse{是的。当对作用量在全时空上积分时，边界条件是已理解的。区别在于，普通时空积分可以通过求和的极限来理解，而Grassmann积分则通过其在有限基\(1,\theta\)上的作用以代数方式定义。}
\end{classroomqa}

\subsection{一个变量}

一个奇变量的任意函数都可以写成其系数位于固定的右侧顺序中的形式
\begin{equation}
f(\theta)=a+\theta b.
\end{equation}
由于\(\partial_\theta\theta=1\)，全导数规则意味着
\begin{equation}
\int d\theta\,1=0.
\end{equation}
剩余的归一化被选为
\begin{equation}
\int d\theta\,\theta=1.
\end{equation}
因此，
\begin{equation}
\int d\theta\,f(\theta)=b.
\end{equation}
积分选出最高可用幂次的系数。在本文使用的左导数约定下，
\begin{equation}
\partial_\theta f=b
=\int d\theta\,f.
\end{equation}

\begin{classroomqa}
\audiencequestion{是否存在满足\(\partial_\theta F=\theta\)的函数\(F(\theta)\)？\lecturetimestamp{01:30:05}}
\lecturerresponse{不存在。这样的函数需要有一个正比于\(\theta^2\)的项，但\(\theta^2=0\)。这是另一种看出为什么不存在不定原函数的Grassmann类比的方式。}
\end{classroomqa}

\subsection{两个变量与测度的顺序}

对于一对复变量\(\theta,\bar\theta\)，将一般偶函数写成有序形式
\begin{equation}
f(\theta,\bar\theta)
=a+\bar\theta\psi+\bar\psi\theta+b\,\bar\theta\theta.
\end{equation}
不存在更高阶的项，因为
\begin{equation}
\theta^2=\bar\theta^2=0,
\qquad
\bar\theta\theta=-\theta\bar\theta.
\end{equation}

我们现在固定约定
\begin{equation}
\int d\theta\,d\bar\theta\,
\bar\theta\theta=1.
\end{equation}
在迭代积分中，相邻的\(d\bar\theta\)运算先作用。于是表格为
\begin{align}
\int d\theta\,d\bar\theta\,1&=0,\\
\int d\theta\,d\bar\theta\,\theta&=0,\\
\int d\theta\,d\bar\theta\,\bar\theta&=0,\\
\int d\theta\,d\bar\theta\,\bar\theta\theta&=1.
\end{align}
因此
\begin{equation}
\int d\theta\,d\bar\theta\,
f(\theta,\bar\theta)=b.
\end{equation}

\begin{classroomqa}
\audiencequestion{对于正比于\(\theta\)的项，真的是\(\theta\)积分使二重积分消失吗？\lecturetimestamp{01:34:59}}
\lecturerresponse{不；实际解释中角色是反过来的。将\(\theta\)对\(\theta\)积分得到一。结果中不再含有\(\bar\theta\)，而剩余的对一的\(\bar\theta\)积分消失。无论哪种方式，只有当两个变量都出现时，二重积分才不为零。}
\end{classroomqa}

\begin{classroomqa}
\audiencequestion{如果交换\(\theta\)和\(\bar\theta\)，积分会改变符号吗？\lecturetimestamp{01:36:33}}
\lecturerresponse{是的。奇变量和奇测度反交换。在上述约定下，
\[
\int d\bar\theta\,d\theta\,\bar\theta\theta=-1.
\]
因此必须选择一种顺序并一致地保持它。物理与约定无关，但中间符号并非如此。}
\end{classroomqa}

对于\(n\)个独立的Grassmann变量，模式是直接的：
\begin{equation}
\int d\theta_n\cdots d\theta_1\,
f(\theta_1,\ldots,\theta_n)
\end{equation}
提取有序最高单项式
\(\theta_1\cdots\theta_n\)的系数，直到由所选变量和测度顺序固定的符号。这一简单的分量提取规则是超空间作用量背后的引擎。超对称变分成为一个奇全导数，而超空间积分将其丢弃，正如普通作用量在适当边界条件下丢弃时空全导数一样。

\section{为什么如此小的微积分有如此大的威力}

这些规则看起来几乎微不足道：
\begin{equation}
\int d\theta\,1=0,
\qquad
\int d\theta\,\theta=1.
\end{equation}
然而它们让我们能够构建相互作用场论，其中玻色子与费米子质量和耦合之间的关系自动保持。超空间微积分的目的与相对论中张量微积分的目的相同：安排记号，使所需的对称性在每一次允许的操作下都得以保持。

符号记账就是代价。每交换两个奇对象就贡献一个负号，而一个计算在概念上可能正确，却会因一次排序选择而出错。这一困难并非装饰性的；它是费米子统计的代数影子。

\begin{classroomqa}
\audiencequestion{格拉斯曼数是与超对称一起被发明的，还是更早就已经存在？\lecturetimestamp{01:41:19}}
\lecturerresponse{它们要古老得多。赫尔曼·格拉斯曼在十九世纪发展了其背后的外代数。微分形式已经在使用反交换的基元素；二十世纪的量子场论和超对称则为同一套数学找到了新的物理角色。}
\end{classroomqa}

现在我们已准备好进行具体构造。超场把完整的玻色子-费米子多重态收集在一起，超空间导数保持代数结构，而 Berezin 积分则提取出分量作用量。接下来的任务是在最简单的一维模型中使用这套机制，然后把它带入场论。
